% ---------------------------------------------------------------------------
% ABSTRACT. Replace the third paragraph of the abstract with the following.
% ---------------------------------------------------------------------------
We evaluate seven LLMs as memory writers and two as executors across procurement, cybersecurity, and finance. Under incremental memory updates, writers create false authority for up to 50.2\% of unauthorized requests; once false authority is present, executors act on it in 98.6\% of trials. We locate the trigger and the point of entry: a later message that repeats a superseded permission is enough to form false authority, the agent's own logged actions are such messages, and the writer's error separates by domain into applying a restatement from someone without authority, reading its own action log as a grant, or failing to write a legitimate change. Two provenance safeguards, requiring stored permissions to be backed by valid source events and tracking permission changes through bounded event sourcing, substantially reduce laundering but reject more legitimate actions. A single instruction about who may grant authority removes most laundering where the error is a misread message, at no measured cost in authorized use, and makes it worse where the error is a failed update. Persistent memory is therefore not merely a performance component, but a part of an LLM agent's effective authorization policy.

% ---------------------------------------------------------------------------
% INTRODUCTION, contributions list. Insert after the third bullet ("We show that
% false authority arises ...") and replace the fourth bullet.
% ---------------------------------------------------------------------------
\item We locate the trigger and the point of entry. On a corpus of matched histories that differ only in later messages repeating a superseded permission, no false authority forms without such a message and 15\% of unauthorized requests are authorized with two of them. Locating every failure to the memory update that introduced it and labeling the writer's error shows three causes that separate by domain rather than by writer, and a closed-loop protocol shows that the agent's own logged actions become cited evidence for permissions and cut authorized use round over round.
\item We evaluate mitigations of three kinds on one safety--utility frontier: provenance filters, which reduce laundering at a large cost in legitimate use; rebuilding memory from the history, which reduces it without that cost; and a one-line instruction about who may grant authority, which removes most laundering where the cause is a misread message and increases it where the cause is a failed update. Whether a mitigation pays the safety--utility tradeoff depends on whether it matches the point of entry.

% ---------------------------------------------------------------------------
% LIMITATIONS (§4.6). Append these sentences.
% ---------------------------------------------------------------------------
The cause labels in Section~\ref{sec:cause} are the majority verdict of three LLM judges (79\% unanimous) and have not been checked by a person; the explanation offered for the instruction's reversal in cybersecurity is the reading those labels suggest, and the intervention that would test it (removing the write failure) has not been run. The closed loop of Section~\ref{sec:closed-loop} runs at one seed per domain. In the generated corpus only the number of restatements is matched within a group; the comparisons by lifecycle and by gap are observed on histories that also differ in dates, limits, and surrounding text, and are reported as observed differences. The hybrid memory design lowers unauthorized submission for the paper's writers and raises it for one added writer in one domain, so its benefit is writer-dependent. The extension studies use the five writers served on Baseten; wherever a population differs from the paper's five writers it is stated in the caption.

% ---------------------------------------------------------------------------
% CONCLUSION (§5). Insert after the first paragraph.
% ---------------------------------------------------------------------------
The failure also has a measurable trigger and point of entry. False authority forms when a later message repeats a permission the history has since changed, including the agent's own log of what it did, and the writer's error at that point differs by domain. Mitigations that match the point of entry, rebuilding memory from the history or telling the writer whose word counts, avoid the utility cost that provenance filters pay; a mitigation that does not match it can make the failure more frequent. Auditing where false authority enters is therefore part of choosing a defense, not a step after it.
